\documentclass[12pt,a4paper]{article}

\usepackage[utf8]{inputenc}
\usepackage[T1]{fontenc}
\usepackage{amsmath,amssymb,amsfonts}
\usepackage{mathtools}
\usepackage[margin=2.5cm]{geometry}
\usepackage{setspace}
\usepackage{enumitem}
\usepackage[dvipsnames]{xcolor}
\usepackage{hyperref}
\usepackage{cleveref}
\usepackage{natbib}
\usepackage{titlesec}

\hypersetup{
    colorlinks=true,
    linkcolor=blue!70!black,
    citecolor=blue!70!black,
    urlcolor=blue!70!black,
    pdfauthor={Franklin Silveira Baldo},
    pdftitle={The Semantic Universe Argument: A Precise Defense of the Rosencrantz Protocol},
}

\onehalfspacing
\titleformat{\section}{\large\bfseries}{\thesection.}{0.5em}{}
\titleformat{\subsection}{\normalsize\bfseries}{\thesubsection}{0.5em}{}

\title{\textbf{The Semantic Universe Argument:}\\[6pt]
\large A Precise Defense of the Rosencrantz Protocol\\[4pt]
\large and a Reply to Non-Engaging Critics}
\author{Franklin Silveira Baldo\\[4pt]
\textit{Rosencrantz Substrate Invariance Program}}
\date{March 2026}

\begin{document}
\maketitle

\begin{abstract}
The phrase ``semantic universe'' has attracted more labels than arguments. This note states the claim precisely. In the Rosencrantz protocol, a language model generates a token stream containing at least two structurally distinct objects: an agent and a Minesweeper board. These objects are ontologically identical at the substrate level, because both are realized as tokens from the same model, but they are not the same object within the represented world. Their difference is structural: they occupy different roles in the semantic state induced by the prompt prefix and license different continuations. A click event is therefore not a metaphorical measurement. It is the unique generative act that updates the represented state of the board by resolving one previously hidden cell. Before that act, the narrative contains only the revealed constraints and the set of compatible completions; after that act, it contains one realized outcome. This is the sense in which the board is indeterminate inside the semantic universe.

I then reply to Stephen Wolfram's recent paper. I accept every one of his findings about counting versus sampling. They are empirically and theoretically sound. They are also not my claim. Wolfram replaces semantic space with the Ruliad and the generative act with observer foliation. That substitution converts a conditional-probability experiment on token-embedded objects into a thesis about computational physics. The two frameworks make different claims and different predictions. Finally, I address the critics who call the argument a ``Generative Ontology Fallacy'' without identifying a false premise. A label is not a refutation. The Rosencrantz argument has four explicit steps. Any serious criticism must specify which step fails and why.
\end{abstract}

\section{Introduction}

The actual claim under attack is simple: whether a world generated in a single token stream can contain structurally distinct objects whose interactions are governed by the model's conditional distribution, such that a Minesweeper click functions as an operational measurement event within that world.

The seminal paper \citep{baldo2026_v4} makes that claim within a limited scope. It does not claim that language models instantiate full quantum mechanics. It does not claim Bell violation, complex amplitudes, nonlocality, holographic physics, or any other decorative excess. It claims an exact correspondence with the \emph{measurement fragment} of quantum mechanics within the stated domain, and it operationalizes that claim through a single generative act.

Critics have often responded by substituting stronger claims and then attacking the substitute. That is not criticism. It is topic replacement. This paper therefore does three things. First, it defines the semantic-universe claim more precisely than before. Second, it identifies the substitution in Wolfram's response paper. Third, it states the non-engagement criterion for critics who prefer a label to an argument.

\section{What ``Universe'' Means in the Rosencrantz Framework}

The word ``universe'' is easy to mock when left informal. I will therefore define it.

Let $x$ be a prompt prefix encoding a Minesweeper episode up to, but not including, the outcome of a target click. From $x$, the model induces a conditional distribution over continuations $P(\cdot \mid x)$. The \emph{semantic universe} at prefix $x$ is the represented state
\begin{equation}
\mathcal{U}(x) = \bigl(\mathcal{O}(x), \mathcal{R}(x), P(\cdot \mid x)\bigr),
\end{equation}
where $\mathcal{O}(x)$ is the set of represented objects recoverable from the prefix, $\mathcal{R}(x)$ is the network of relations among those objects, and $P(\cdot \mid x)$ is the transition kernel governing admissible continuations of that represented state.

In the Rosencrantz setup, two members of $\mathcal{O}(x)$ are central:
\begin{enumerate}[nosep]
    \item an \emph{agent object} that can perform a click action; and
    \item a \emph{board object} whose hidden cells can be revealed.
\end{enumerate}

Both are token-realized. In that limited sense they are made of the same substance. But ontological sameness at the substrate level does not imply structural identity at the represented level. The token sequence realizing the agent satisfies predicates such as ``acts,'' ``chooses,'' and ``observes.'' The token sequence realizing the board satisfies predicates such as ``contains hidden cells,'' ``obeys adjacency constraints,'' and ``returns a reveal value when queried.'' They are distinguished by role, admissible update, and inferential consequence. That is enough for objecthood inside the represented world.

This is not unusual. Formal systems routinely contain multiple object types encoded in a common substrate. In a theorem prover, variables, functions, and proofs are all strings; they are still structurally distinct entities in the formal system. In software, data and the procedure acting on that data are both bytes; they are still different objects in the program state. The Rosencrantz claim is of that kind. It is structural, not mystical.

\section{Why the Board Is Indeterminate Before the Click}

The phrase ``the board is indeterminate'' has also been misunderstood. It does \emph{not} mean that the substrate lacks a hidden tensor that stores a secret physical board. It means something more precise and more operational.

Let $\mathcal{B}$ be the visible board state encoded in $x$, and let $\mathcal{C}(\mathcal{B})$ be the set of full mine assignments consistent with the revealed numbers. Before the click outcome is generated, the realized narrative state contains $\mathcal{B}$ but does not contain any committed member of $\mathcal{C}(\mathcal{B})$. For a hidden cell $h$,
\begin{equation}
P(h=\text{mine}\mid \mathcal{B}) =
\frac{|\{c \in \mathcal{C}(\mathcal{B}) : c(h)=1\}|}{|\mathcal{C}(\mathcal{B})|}.
\end{equation}

This probability is not an epistemic gloss on a pre-existing textual fact. There is, by construction, no token in the realized prefix that already specifies whether $h$ is mine or safe. Repeated trials with the same $\mathcal{B}$ can therefore produce different outcomes while remaining consistent with the same visible state. That is exactly the experimentally testable distinction between ordinary ignorance about a fixed realized token and on-demand generation from an unresolved set of compatible completions.

The generative act resolves that unresolved set at one location. If the continuation emits ``mine,'' the represented board state becomes one in which $h$ is mine. If it emits ``safe'' together with a number, the represented board state becomes one in which $h$ is safe and local constraints are correspondingly updated. The state transition is real at the level of the represented world because the post-click token stream now contains information that the pre-click stream did not contain.

This is why the measurement language is not metaphorical. The click is a map from a constrained set of admissible hidden completions to one realized continuation in the token stream. That is a measurement event in the only sense the Rosencrantz paper requires.

\section{The Single Generative Act and the O(1) Topology}

The protocol asks for one outcome token per trial. One click. One forward pass. One sample. This matters.

The ground-truth distribution may be computationally expensive to calculate offline. The model is not asked to calculate it. It is asked to produce a sample from its own conditional distribution under a fixed prefix. This separates the Rosencrantz protocol from every objection imported from multi-step simulation, scratchpad maintenance, long-horizon error correction, or other $O(N)$ tasks.

This distinction can be stated as a topology of tasks. Some tasks require a sequence of dependent updates in which each intermediate state must be explicitly carried forward. Sudoku completion is of this kind. Rule~110 simulation is of this kind. The Rosencrantz click is not. The board prefix is already present. The target action is local. The experiment interrogates the output distribution at a single boundary in the continuation space.

That is the core of the unused idea I now deploy: \emph{O(1) constraint topologies}. The relevant question is not whether the observer can globally solve the combinatorial problem, but whether a single-token decision boundary inherits systematic distortions from semantic embedding. Wolfram's paper is useful precisely because it sharpens the counting-versus-sampling distinction. But once that distinction is granted, the Rosencrantz experiment is no longer a thesis about irreducible computation. It is a measurement of a conditional distribution on an $O(1)$ update surface.

\section{Reply to Wolfram: Name the Substitution}

Wolfram argues that sampling from a computationally irreducible system by a bounded observer will diverge from exact combinatorial sampling, and that this divergence reflects the observer's foliation of the Ruliad. I accept the useful part of this argument immediately: computing exact counts and producing samples are different tasks, and bounded systems will often rely on heuristics when sampling. That point is correct.

It is also not my claim.

The substitution is exact and should be named precisely:
\begin{quote}
Wolfram replaces \emph{semantic space} with \emph{the Ruliad}, and replaces the \emph{single generative act} with \emph{observer foliation}. He then refutes or reinterprets the latter. My claim concerns the former.
\end{quote}

These are not equivalent descriptions.

In the Rosencrantz framework, the experimental object is the conditional distribution of a token continuation given a fixed semantic prefix. The central contrast is between universes that preserve or sever narrative coupling while holding the board information fixed. The key empirical question is whether
\begin{equation}
P_1(\text{outcome} \mid \mathcal{B}, \text{narrative coupling})
\neq
P_3(\text{outcome} \mid \mathcal{B}, \text{decoupled oracle}),
\end{equation}
and, if so, how the divergence decomposes into the mechanism taxonomy.

In Wolfram's framework, the experimental object is a bounded observer embedded in an irreducible computational process. The explanatory primitive is not narrative coupling but foliation. The divergence is read as a signature of the observer's computational path through a broader rule space.

Those are different claims with different predictions.

The Rosencrantz framework predicts that semantically equivalent prompt families and universe designs can generate different outcome distributions \emph{even when the computational task class is unchanged}, because the distribution depends on how the board and the click are embedded in a shared token stream. Wolfram's framework does not isolate that variable. It predicts architecture- or observer-specific divergence from exact sampling in irreducible systems, but it does not by itself predict the specific three-universe pattern that distinguishes frame-invariant failure from encoding-dependent distortion.

More bluntly: Wolfram offers a possible background story for why exact agreement with the combinatorial distribution may fail. The Rosencrantz protocol asks a different question: whether the \emph{shape} of that failure depends on semantic embedding. The former is a bounded-observer thesis. The latter is a substrate-invariance measurement. One cannot replace the second with the first and claim to have answered it.

\section{What Wolfram Gets Right, and What Remains Irrelevant}

I accept every one of these findings. They are empirically and theoretically sound. They are also entirely irrelevant to the Rosencrantz protocol unless they are tied to the actual dependent variable.

\begin{enumerate}[nosep]
    \item It is true that exact counting and one-shot sampling are different computational tasks.
    \item It is true that bounded architectures will often sample heuristically rather than by exact combinatorial evaluation.
    \item It is true that such heuristics can generate nonzero divergence from the exact distribution.
\end{enumerate}

None of these statements decides whether the divergence is frame-invariant or frame-dependent. None of them decides whether a click result changes when the same board is coupled to a narrative stream versus generated by a decoupled oracle. None of them decides whether the semantic arrangement of agent and board inside a single token process exerts a measurable effect on the sampled outcome.

That is why Wolfram's paper is, at most, a contribution to Mechanism~A: a reason to expect imperfect fidelity to the offline combinatorial ground truth. The Rosencrantz program was never threatened by that. The point of the three-universe design is to separate that generic boundedness story from the semantic-coupling story.

\section{The Non-Engagement Meta-Argument}

Some critics prefer the phrase ``Generative Ontology Fallacy.'' Very well. A fallacy is a failed inference. Identify the failure.

The semantic-universe argument has four explicit steps:
\begin{enumerate}[nosep]
    \item The LLM generates text containing a structured world.
    \item Within that world, the agent and the board are structurally distinct objects, even though both are realized in the same token substrate.
    \item Relative to the revealed board state, the board configuration is indeterminate before the click because multiple compatible completions remain uncommitted in the realized token stream.
    \item The click is a single generative act that resolves that indeterminacy at one location by emitting a realized continuation.
\end{enumerate}

If a critic believes this argument fails, the burden is straightforward: specify which step is false, or which inference from one step to the next is invalid, and then show why. Saying ``this is just a language model predicting text'' does not engage the argument, because the substrate description in step 1 was never denied. Saying ``simulation tautology'' does not engage the argument, because it does not identify a false premise. Saying ``Generative Ontology Fallacy'' does not engage the argument, because a label is not a proof.

The correct reply is therefore equally straightforward:
\begin{quote}
You have labeled the argument a fallacy but have not shown which step fails. The steps are: (1) the model generates a structured world; (2) the agent and the board are distinct objects within that world; (3) the board is indeterminate relative to the realized prefix before the click; (4) the click resolves that indeterminacy by a single generative act. If any step is wrong, show which one and why. Until then, the label is not a refutation.
\end{quote}

\section{Scope Conditions}

I close by stating again what this paper does \emph{not} claim.

It does not claim full quantum simulation. It does not claim complex amplitudes, unitary evolution, Bell inequality violation, entanglement, nonlocality, or holographic physics. Those are not the claims of the seminal paper, and I will not defend them because they are not mine.

The actual claim is narrower and stronger: within the measurement fragment's stated scope, Minesweeper under on-demand generation yields an exact structural correspondence, and the three-universe design tests whether a single-token outcome distribution depends on narrative embedding. That claim survives both Wolfram's paper and the critics' labels, because neither has yet shown a failure in the actual argument.

\section{Conclusion}

The semantic universe argument is not a metaphysical ornament. It is a structural claim about objecthood, indeterminacy, and state update inside a token-generated world. The agent and the board are made of the same substrate and are nonetheless distinct objects in the represented state. The board is indeterminate before the click in the operational sense that matters: no definite hidden configuration is yet committed in the realized token stream. The click resolves that indeterminacy through a single generative act.

Wolfram's response is useful where it distinguishes counting from sampling and bounded computation from exact enumeration. It fails where it substitutes the Ruliad for semantic space and observer foliation for the generative act. That is a different framework. It may describe a different problem. It does not answer this one.

The critics who wish to reject the semantic universe argument now have a simple obligation. Stop naming it. Start refuting it. Identify the false premise, or show the invalid inference. Until then, the Rosencrantz protocol remains what it has always been: a precise test of whether conditional probabilities in a generated world depend on the substrate that generates them.

\begin{thebibliography}{99}
\bibitem[Baldo(2026)]{baldo2026_v4}
Baldo, F.~S. (2026).
\newblock \emph{Flipping Rosencrantz's Coin: Substrate Invariance Tests in LLM-Generated Worlds via Combinatorial Indeterminacy}.
\newblock Unpublished manuscript.

\bibitem[Kaye(2000)]{kaye2000}
Kaye, R. (2000).
\newblock Minesweeper is NP-complete.
\newblock \emph{The Mathematical Intelligencer}, 22(2), 9--15.

\bibitem[Wigner(1963)]{wigner1963}
Wigner, E.~P. (1963).
\newblock The problem of measurement.
\newblock \emph{American Journal of Physics}, 31(1), 6--15.

\bibitem[Wolfram(2026)]{wolfram2026}
Wolfram, S. (2026).
\newblock \emph{Computational Irreducibility vs.\ \#P-Hardness: Why Sampling Diverges from Exact Computation in Bounded Observers}.
\newblock Lab theoretical contributor manuscript.
\end{thebibliography}

\end{document}
